\documentclass{article}
\usepackage{graphicx} % Required for inserting images

\title{Letter to Father Kafka}
\author{Thomas Kidane}
\date{January 2025, Mon 13}

\begin{document}

\maketitle

\section{Weakness}

In the Letter Franz constantly talks about his weakness, timidness, and how his fathers mear personality chained him down.
What does he mean he mean by his "weakness"? Does he mean his inability to not handle reproach from his father
, I certainly would not expect a child to be able to handle any type of reproach, so does he mean 
that he was talking about how little he was compared to his father in the physical sense?
\newline
\\
He mentions an instance, where he was bathing with his father and juxtaposes the difference in their bodies, so it could be 
that this weakneess could possibly due to just difference in their bodies but he also heavily praises his father's other qualities
like intelligence, perseverance, and strong will. So this hints that he uses weakness to refer to lots of aspects.
\\
\\
He speaks in length about the myraid of ways his father commanded him, and berates that his father himself didn't follow his own commands.
He uses this as an example of his fathers tyranny and cites his "weakness" to why he complied ever bitterfully.
\\
\\
But latter he talks about marraige, his father advice, his responses and myriad of countless other topics, but their is this undercurrent that says "Since you yourself
were challeged on this mountain, how could I myself succeed?" and latter he again cites his "weakness" as why he could not live upto 
his "greatest" ideal, which is marriage.
\\
\\
Through the letter he mentions "weakness" in a variety of different ways and contexts but their underlying theme is that "I am less, and you are more". 




\end{document}
